\section{Time Hierarchy --- Clock Loop}
%%% generated by the blueprint pipeline from TCSlib.Complexity.TimeHierarchy.ClockLoop
\subsection*{Time Hierarchy --- Clock Loop: Introduction}

Declarations of the Lean module \texttt{TCSlib.Complexity.TimeHierarchy.ClockLoop}.

\subsection*{Time Hierarchy --- Clock Loop: Results}

\begin{lemma}[Updating the middle symbol on a buffer tape]
\lean{Complexity.TimeHierarchy.bufferTape_update_mid}
\label{Complexity.TimeHierarchy.bufferTape_update_mid}
\leanok
\uses{Turing.FinTM.bufferTape}
\difficulty{4}
For finite Boolean words $l$ and $m$ and Boolean symbols $a$ and $b$, updating the tape
storing the concatenation of $l$, $a$, and $m$ at the integer position $|l|$ to contain
$b$ yields exactly the tape storing the concatenation of $l$, $b$, and $m$.
\end{lemma}

\begin{lemma}[Borrowing steps in the clock computation]
\lean{Complexity.TimeHierarchy.borrow_run}
\label{Complexity.TimeHierarchy.borrow_run}
\leanok
\uses{Complexity.TimeHierarchy.OutReg, Complexity.TimeHierarchy.bufferTape_update_mid, Complexity.TimeHierarchy.clockTM, Complexity.TimeHierarchy.clockTM_step_some, Complexity.TimeHierarchy.clockTr, Complexity.TimeHierarchy.ctrAction_apply, Complexity.TimeHierarchy.ctrCfg, Complexity.TimeHierarchy.ctrCfg_ctr, Turing.Cfg, Turing.FinTM.bufferTape, Turing.FinTM.bufferTape_nat, Turing.MultiTapeTM.runFrom, Turing.MultiTapeTM.runFrom_succ_eq_step, Turing.MultiTapeTM.step}
\difficulty{4}
Let $K$ and $W$ be finite Turing machines over the Boolean alphabet, let $x$ be a finite
Boolean input, let $c$ be a configuration of $W$ on $x$, let $q$ be a state of $W$, let
$r$ be an output register, let $i,j\in\mathbb{N}$, let $\mathrm{rest}$ be a finite
Boolean word, and let $kt$ and $kh$ specify the contents and head positions of the work
tapes of $K$. Start the clock machine for $K$ and $W$ with simulated configuration $c$,
simulated state $q$, register $r$, counter-tape contents storing $i$ true symbols
followed by $j$ false symbols and then $\mathrm{rest}$, counter head at position $i$,
the specified tapes of $K$, and empty output. After $j$ steps, its state is still the
decrement state for $q,r$, its simulated configuration and $K$ tapes are unchanged, and
the counter tape stores $i+j$ true symbols followed by $\mathrm{rest}$ with its head at
position $i+j$; its output remains empty.
\end{lemma}

\begin{lemma}[Returning from the clock scan]
\lean{Complexity.TimeHierarchy.ret_run}
\label{Complexity.TimeHierarchy.ret_run}
\leanok
\uses{Complexity.TimeHierarchy.OutReg, Complexity.TimeHierarchy.clockTM, Complexity.TimeHierarchy.clockTM_step_some, Complexity.TimeHierarchy.clockTr, Complexity.TimeHierarchy.ctrAction_apply, Complexity.TimeHierarchy.ctrCfg, Complexity.TimeHierarchy.ctrCfg_ctr, Turing.Cfg, Turing.MultiTapeTM.runFrom, Turing.MultiTapeTM.runFrom_succ_eq_step, Turing.MultiTapeTM.runFrom_succ_eq_step', Turing.MultiTapeTM.runFrom_zero, Turing.MultiTapeTM.step}
\difficulty{4}
Let $K$ and $W$ be finite Boolean-alphabet Turing machines, and let $c$ be a
configuration of $W$ on Boolean input $x$, with state $q$ and output register $r$.
Choose work-tape contents $kt$ and head positions $kh$ for the tapes corresponding to
$K$, and let $\tau:\mathbb Z\to\mathrm{Option}(\mathrm{Bool})$ describe the clock tape.
Suppose $\tau(-1)$ is blank. For every natural number $j$ such that $\tau(i)$ is
nonblank for every natural number $i<j$, running the clock machine for $j+1$ steps from
the configuration with state $\mathrm{ret}(q,r)$, underlying configuration $c$, clock
tape $\tau$, clock-head position $j-1$, tapes $kt,kh$, and empty output yields the
configuration with state $\mathrm{dec}(q,r)$, the same underlying configuration and
tapes, clock tape $\tau$, clock-head position $0$, and empty output.
\end{lemma}

\begin{lemma}[Clock simulation step]
\lean{Complexity.TimeHierarchy.wstep}
\label{Complexity.TimeHierarchy.wstep}
\leanok
\uses{Complexity.TimeHierarchy.OutReg, Complexity.TimeHierarchy.OutReg.push, Complexity.TimeHierarchy.clockTM, Complexity.TimeHierarchy.clockTM_step_some, Complexity.TimeHierarchy.clockTr, Complexity.TimeHierarchy.ctrCfg, Complexity.TimeHierarchy.ctrCfg_ctr, Complexity.TimeHierarchy.ctrCfg_right, Turing.Action.apply, Turing.Cfg, Turing.Cfg.inputSymbol, Turing.Cfg.workTapeSymbols, Turing.MultiTapeTM.step}
\difficulty{5}
Let $K$ and $W$ be finite Turing machines over the Boolean alphabet, let $x$ be an input
bit string, and let $c$ be a configuration of $W$ on $x$ whose state is $q$. Let $r$ be
an output register, let $\tau$ be a tape content function with $\tau(z)=1$, and let $z$
be an integer; for any fixed contents and head positions of the other tapes, one step of
the clock machine from the configuration encoding $c$, with clock state indicating that
it is simulating $W$ in its decoding phase, equals the encoded configuration after one
step of $W$: the simulated state is updated according to $W$'s transition, the clock
cell at $z$ is changed to blank and its head moves to $z-1$, and the output is the
one-bit value indicating whether the updated register is not exactly the register
containing only $1$ if $W$'s transition halts, and is empty otherwise.
\end{lemma}

\begin{lemma}[Clock underflow halts]
\lean{Complexity.TimeHierarchy.underflow_step}
\label{Complexity.TimeHierarchy.underflow_step}
\leanok
\uses{Complexity.TimeHierarchy.OutReg, Complexity.TimeHierarchy.clockTM, Complexity.TimeHierarchy.clockTM_step_some, Complexity.TimeHierarchy.clockTr, Complexity.TimeHierarchy.ctrCfg, Complexity.TimeHierarchy.ctrCfg_ctr, Turing.Action.apply, Turing.Cfg, Turing.MultiTapeTM.step}
\difficulty{3}
Let $K$ and $W$ be finite Boolean-alphabet Turing machines, let $c$ be a configuration
of $W$ on Boolean input $x$ with state $q$ and output register $r$, and choose contents
$kt$ and head positions $kh$ for the tapes corresponding to $K$. Let $\tau$ be a
clock-tape content function and $z$ an integer. If the clock tape is blank at position
$z$, then one step of the clock machine from the configuration with state
$\mathrm{dec}(q,r)$, underlying configuration $c$, clock tape $\tau$, clock-head
position $z$, tapes $kt,kh$, and empty output halts and has output consisting of the
single bit $1$.
\end{lemma}

\begin{definition}[Counter bound]
\lean{Complexity.TimeHierarchy.ctrBound}
\label{Complexity.TimeHierarchy.ctrBound}
\leanok
\uses{Complexity.TimeHierarchy.ctrPop, Complexity.TimeHierarchy.ctrVal}
For finite Turing machines $K$ and $W$ over the Boolean alphabet and a Boolean string
$s$, define the counter bound by $4c(s)+2(|s|-p(s))+|s|+1$, where $c(s)$ denotes the
counter value and $p(s)$ the counter population count.
\end{definition}

\begin{definition}[Loop answer predicate]
\lean{Complexity.TimeHierarchy.loopAnswer}
\label{Complexity.TimeHierarchy.loopAnswer}
\leanok
\uses{Turing.Cfg, Turing.MultiTapeTM.runFrom}
Let $W$ be a finite Turing machine over the Boolean alphabet, let $x$ be a Boolean input
string, let $c$ be a configuration of $W$ on input $x$, and let $v$ be a natural number.
The loop answer is true exactly when running $W$ from $c$ for $v$ steps does not both
halt and produce the one-symbol output string consisting of true.
\end{definition}

\begin{theorem}[Clocked simulation of a bounded computation]
\lean{Complexity.TimeHierarchy.clockTM_loop}
\label{Complexity.TimeHierarchy.clockTM_loop}
\leanok
\uses{Complexity.TimeHierarchy.OutReg.ofList, Complexity.TimeHierarchy.OutReg.ofList_eq_one_true, Complexity.TimeHierarchy.OutReg.ofList_push, Complexity.TimeHierarchy.OutReg.push, Complexity.TimeHierarchy.borrow_run, Complexity.TimeHierarchy.bufferTape_update_mid, Complexity.TimeHierarchy.clockTM, Complexity.TimeHierarchy.ctrBound, Complexity.TimeHierarchy.ctrCfg, Complexity.TimeHierarchy.ctrPop, Complexity.TimeHierarchy.ctrPop_le_length, Complexity.TimeHierarchy.ctrPop_replicate, Complexity.TimeHierarchy.ctrVal, Complexity.TimeHierarchy.ctrVal_borrow, Complexity.TimeHierarchy.ctrVal_pos_decomp, Complexity.TimeHierarchy.ctrVal_zero_eq, Complexity.TimeHierarchy.loopAnswer, Complexity.TimeHierarchy.ret_run, Complexity.TimeHierarchy.underflow_step, Complexity.TimeHierarchy.wstep, Turing.Action.apply, Turing.Cfg, Turing.Cfg.inputSymbol, Turing.Cfg.workTapeSymbols, Turing.FinTM.bufferTape, Turing.FinTM.bufferTape_left, Turing.FinTM.bufferTape_nat, Turing.MultiTapeTM.runFrom, Turing.MultiTapeTM.runFrom_add, Turing.MultiTapeTM.runFrom_of_halt, Turing.MultiTapeTM.runFrom_succ_eq_step, Turing.MultiTapeTM.runFrom_succ_eq_step', Turing.MultiTapeTM.step}
\difficulty{6}
Let $K$ and $W$ be finite Turing machines over the Boolean alphabet, let $x$ and $s$ be
Boolean words, and let $c$ be a configuration of $W$ on input $x$. If $c$ is in state
$q$, then for every assignment of contents and head positions to the work tapes of $K$,
the clock machine started with simulated configuration $c$, an output register recording
the contents of $c$’s output tape, and counter tape containing $s$ halts within
$4v+2(|s|-r)+|s|+1$ steps and outputs the one-bit word whose bit is true exactly when it
is not the case that after $v$ steps from $c$, $W$ has halted and its output is the
one-bit word true; here $v$ is the integer represented by $s$ in
least-significant-bit-first order and $r$ is the number of true bits in $s$.
\end{theorem}

\begin{theorem}[Clock machine diagonal specification]
\lean{Complexity.TimeHierarchy.clockTM_spec}
\label{Complexity.TimeHierarchy.clockTM_spec}
\leanok
\uses{Complexity.TimeHierarchy.OutReg.ofList, Complexity.TimeHierarchy.clockTM, Complexity.TimeHierarchy.clockTM_loop, Complexity.TimeHierarchy.clockTM_setup, Complexity.TimeHierarchy.ctrBound, Complexity.TimeHierarchy.ctrVal, Complexity.TimeHierarchy.loopAnswer, Turing.FinTM.ComputesInTime, Turing.FinTM.ComputesInTime.mono, Turing.FinTM.computesInTime_iff, Turing.MultiTapeTM.initCfg, Turing.MultiTapeTM.runFrom, Turing.MultiTapeTM.runFrom_add}
\difficulty{5}
Let $K$ and $W$ be finite Turing machines over the Boolean alphabet, and let $x$ and $s$
be finite Boolean words. If $K$ halts on input $x$ with output $s$ after exactly $t_K$
steps, then the clock machine built from $K$ and $W$ halts on input $x$, outputs the
one-symbol word whose value is the negation of the proposition that $W$ halts on $x$
within exactly the number of steps represented by $s$ and outputs the one-symbol word
containing true, and does so after exactly $t_K+4|s|+|x|+4\,\operatorname{ctrVal}(s)+6$
steps.
\end{theorem}
